\documentclass[10pt,a4paper]{amsart}
\usepackage[margin=19mm]{geometry}
\usepackage{amsmath,amssymb,mathtools}
\usepackage[scr=rsfs,cal=euler]{mathalfa}
\usepackage{xfrac}
\usepackage[unicode,hidelinks,bookmarks=false]{hyperref}
\newcommand{\ie}{i.e.\ }
\newcommand{\cf}{cf.\ }
\newcommand{\resp}{respectively\ }
\newcommand{\Subsection}{\subsection}
\providecommand{\nobreakdash}{\nobreak\hbox{-}}
\setlength{\emergencystretch}{2em}
\multlinegap0pt
\usepackage{tikz}
\usetikzlibrary{cd}
\usepackage{amssymb}
\usepackage{bbm}
\usepackage{enumerate}
\usepackage[all]{xy}
\newtheorem{proposition}{Proposition}
\newcommand{\TwSp}{\mathrm{TwSp}}
\newcommand{\bS}{\mathbb{S}}
\newcommand{\och}{\quad \text{and} \quad}
\title{Twisted Spectra Revisited}
\author{Alice Hedenlund, Tasos Moulinos}
\date{Source: arXiv:2312.07403v2 (CC BY 4.0)}
\begin{document}
\maketitle

\noindent\emph{Source and licence.} Twisted Spectra Revisited. Version arXiv:2312.07403v2 (CC BY 4.0), distributed by arXiv under the Creative Commons Attribution 4.0 International licence, http://creativecommons.org/licenses/by/4.0/, as recorded in the arXiv licence metadata for this exact version and shown on the abstract page https://arxiv.org/abs/2312.07403v2. Full licence notice: \texttt{licenses/arxiv-2312.07403-CC-BY-4.0.txt}.\par\smallskip

\noindent\textit{Editorial scope.} Counted unit: the article's proposition proposition (source0.tex lines 2518--2523). The article's complete original proof of it is delivered with it (source0.tex lines 2525--2580, one \texttt{proof} environment of 56 lines). No mathematics is written, repaired or re-derived: the statement and the proof are the article's own lines, in the article's own notation, and no retained line is substituted or re-flowed.  No further source-local context is needed: every cross-reference of every retained line resolves inside the two delivered spans.  Nothing was omitted from any retained span, so the package carries no omission record.  No retained line contains a citation, so the wrapper carries no bibliography and no bibliography entry.  The retained lines contain the article's own diagram code, which the wrapper typesets with the standard tikz packages; no image file is delivered and the package declares no asset dependency.  Editorial formatting only, in addition: an AMS \texttt{amsart} preamble that restates the article's own declarations and macros used by the retained lines, namely the environments \texttt{proposition} on separate counters, together with the standard packages those lines need.  The retained environments print in this document as Proposition 1 in order of appearance, whereas the article prints the same items with the numbering of its own full document.  The complete native source of the article as distributed in the arXiv e-print for this exact version is bundled unchanged as \texttt{sources/151389/source0.tex}.
\begin{proposition}
  Let $X \in \TwSp(A \times B , -\tau \boxtimes \sigma)$ and let $Y \in \TwSp(A \times B, -\sigma \boxtimes \rho)$. Then
  \[
  \Phi_{Y \odot X} \simeq \Phi_{Y} \circ \Phi_{X} \och \Phi_{\Delta^{\tau}_! \bS} \simeq 1_{\TwSp(A,\tau)} \,.
  \]
\end{proposition}

\begin{proof}
  The assertions essentially follows from using the projection isomorphism and Beck--Chevalley a couple of times for various maps.
  To keep track of maps let us use the notation
  \begin{equation*}
  \begin{aligned}
  p_A &: A \times B \to A \quad \quad  &q_B: B \times C \to B \quad \quad r_{A} &: A \times C \to C \\
  p_B &: A \times B \to B  \quad \quad &q_C: B\times C \to B \quad \quad  r_C &: B \times C \to C
  \end{aligned}
  \end{equation*}
  for the various projections. We will also write
  \[
  p_{A \times B} : A \times B \times C \to A \times C \quad p_{B \times C} : A \times B \times C \to B \times C \quad p_{A \times C} : A \times B \times C \to A \times C
  \]
  and
  \[
  \pi_{A} : A \times B \times C \to A \quad \pi_{B} : A \times B \times C \to B \quad \pi_{C} : A \times B \times C \to A \,.
  \]
  Per definition, we have that
  \[
  \Phi_Y \circ \Phi_X = (q_C)_! \circ (-\otimes Y) \circ (q_B)^* \circ (p_B)_! \circ (- \otimes X) \circ (p_A)^*  \,,
  \]
  and that

  \begin{align*}
  \Phi_{Y \odot X}& = (r_{C})_! \circ (- \otimes (Y \odot X)) \circ (r_A)^* \\
  &= (r_{C})_! \circ (- \otimes ((p_{A \times C})_!(p^*_{B \times C} Y \otimes p_{A \times B}^* X))) \circ (r_A)^* \,.
\end{align*}

  These are the compositions that we are trying to identify.
  We start by using Beck--Chevalley on the pullback diagram
  \[
  \begin{tikzcd}
    A \times B \times C \arrow[r,"p_{A \times B}"] \arrow[d,"p_{B \times C}"'] & A \times B \arrow[d,"p_B"] \\
    B \times C \arrow[r,"q_B"] & B
  \end{tikzcd}
  \]
  so that we have that $(q_B)^* \circ (p_B)_! \simeq (p_{B \times C})_! \circ (p_{A \times B})^*$. By combining this with the fact that the pullback $(p_{A \times B})^*$ is symmetric monoidal we obtain
  \[
  \Phi_Y \circ \Phi_X \simeq (q_C)_! \circ (-\otimes Y) \circ (p_{B \times C})_! \circ (- \otimes (p_{A \times B})^*(X)) \circ (\pi_A)^* \,.
  \]
  By using the projection isomorphism for the map $p_{B \times C}$ we can further write
  \[
  \Phi_Y \circ \Phi_X \simeq (\pi_C)_! \circ (-\otimes (p_{B \times C})^*(Y)) \circ (- \otimes (p_{A \times B})^*(X)) \circ (\pi_A)^* \,.
  \]
  Let us now look closer at the functor $\Phi_{Y \odot X}$. By using the projection isomorphism for the map $p_{A \times C}$ to make the identification
  \[
  \Phi_{Y \odot X} \simeq (\pi_{C})_! \circ (- \otimes ((p_{B \times C})^* Y \otimes (p_{A \times B})^* X)) \circ (\pi_A)^*
  \]
  which by inspection agrees with the identification we made for $\Phi_Y \circ \Phi_X$. This finishes the proof of the first assertion.
  For the second assertion, we will use the notation $p_1 , p_2 : A \times A \to A$  for the projection onto the first and second factor, respectively.
  Using the projection isomorphism for the diagonal map $\Delta : A \to A \times A$ we obtain
  \begin{align*}
    (p_2)_!\circ (- \otimes \Delta_! \bS)\circ (p_1)^* &\simeq (p_2)_! \circ \Delta_! \circ (- \otimes \bS) \circ \Delta^* (p_1)^* \\ \simeq (-)
  \end{align*}
  where we have used that $p_i \circ \Delta \simeq 1$ for $i=1,2$ and that tensoring with $\bS$ is the identity.
\end{proof}

\end{document}
